% !TeX program = xelatex
% 编译：xelatex YingluoNet_v1_1_structure_manual.tex （运行两次以生成目录）
% 本文件不依赖外部图片。所用结构图均由 TikZ 绘制。
\documentclass[UTF8,11pt,a4paper,fontset=fandol]{ctexart}
\usepackage[left=22mm,right=22mm,top=22mm,bottom=21mm,headheight=16pt,headsep=7mm]{geometry}
\usepackage{amsmath,amssymb,mathtools,bm}
\usepackage{fontspec}
\setmainfont{lmroman10-regular.otf}[BoldFont=lmroman10-bold.otf,ItalicFont=lmroman10-italic.otf,BoldItalicFont=lmroman10-bolditalic.otf]
\setsansfont{lmsans10-regular.otf}[BoldFont=lmsans10-bold.otf,ItalicFont=lmsans10-oblique.otf,BoldItalicFont=lmsans10-boldoblique.otf]
\setmonofont{lmmono10-regular.otf}[Scale=0.86]
\usepackage{graphicx,booktabs,tabularx,array,multirow}
\usepackage[table]{xcolor}
\usepackage{tikz}
\usetikzlibrary{arrows.meta,positioning,calc,fit,backgrounds}
\usepackage[most]{tcolorbox}
\usepackage{fancyhdr,lastpage,setspace,enumitem,tocloft}
\usepackage{hyperref}
\definecolor{Ink}{HTML}{172B41}
\definecolor{Teal}{HTML}{087F8C}
\definecolor{Gold}{HTML}{B87832}
\definecolor{Muted}{HTML}{596A78}
\definecolor{Line}{HTML}{D9E2E7}
\definecolor{Tint}{HTML}{EFF6F7}
\definecolor{Warm}{HTML}{FBF5EC}
\hypersetup{colorlinks=true,linkcolor=Ink,urlcolor=Teal,citecolor=Teal,
  pdftitle={YingluoNet v1.1 结构说明书},pdfauthor={YingluoNet},
  pdfsubject={相位耦合、双分支门控与残差连接的数学结构及前向传播},
  bookmarksnumbered=true,pdfstartview=FitH}
\ctexset{
  section={format=\Large\sffamily\bfseries\color{Ink},beforeskip=0pt,afterskip=13pt},
  subsection={format=\normalsize\sffamily\bfseries\color{Teal},beforeskip=11pt,afterskip=6pt},
  subsubsection={format=\normalsize\sffamily\bfseries,beforeskip=8pt,afterskip=4pt}}
\setstretch{1.19}
\setlength{\parindent}{0pt}
\setlength{\parskip}{6pt}
\setlength{\abovedisplayskip}{8pt}
\setlength{\belowdisplayskip}{8pt}
\setlength{\abovedisplayshortskip}{4pt}
\setlength{\belowdisplayshortskip}{5pt}
\setlength{\emergencystretch}{2em}
\setcounter{tocdepth}{1}
\renewcommand{\cftsecfont}{\normalfont}
\renewcommand{\cftsecpagefont}{\normalfont\color{Muted}}
\setlength{\cftbeforesecskip}{7pt}
\setlength{\cftsecnumwidth}{2.7em}
\renewcommand{\cftsecleader}{\cftdotfill{\cftdotsep}}
\renewcommand{\arraystretch}{1.28}
\newcolumntype{Y}{>{\raggedright\arraybackslash}X}
\newcolumntype{C}[1]{>{\centering\arraybackslash}p{#1}}
\pagestyle{fancy}
\fancyhf{}
\fancyhead[L]{\small\sffamily\color{Muted}YingluoNet\quad /\quad v1.1}
\fancyhead[R]{\small\sffamily\color{Muted}结构说明书}
\fancyfoot[L]{\scriptsize\sffamily\color{Muted}相位耦合\enspace ·\enspace 双分支门控\enspace ·\enspace 残差连接}
\fancyfoot[R]{\small\sffamily\color{Muted}\thepage\ /\ \pageref{LastPage}}
\renewcommand{\headrulewidth}{0.4pt}
\renewcommand{\footrulewidth}{0pt}
\newtcolorbox{explain}[1]{enhanced,colback=Tint,colframe=Tint,boxrule=0pt,
  borderline west={2pt}{0pt}{Teal},arc=1mm,left=10pt,right=10pt,top=8pt,bottom=8pt,
  title={#1},coltitle=Ink,fonttitle=\bfseries\sffamily,
  attach title to upper={\par\smallskip},before skip=9pt,after skip=9pt}
\newtcolorbox{notebox}[1]{enhanced,colback=Warm,colframe=Warm,boxrule=0pt,
  borderline west={2pt}{0pt}{Gold},arc=1mm,left=10pt,right=10pt,top=7pt,bottom=7pt,
  title={#1},coltitle=Ink,fonttitle=\bfseries\sffamily,
  attach title to upper={\par\smallskip},before skip=8pt,after skip=8pt}
\newcommand{\R}{\mathbb R}
\newcommand{\sig}{\operatorname{sigmoid}}
\newcommand{\spplus}{\operatorname{softplus}}
\newcommand{\diag}{\operatorname{diag}}
\newcommand{\YLayer}{\mathcal{Y}}
\newcommand{\purpose}[1]{\begin{explain}{设计目的}#1\end{explain}}
\newcommand{\newsection}[1]{\clearpage\section{#1}}
\tikzset{
  block/.style={draw=Line,fill=white,rounded corners=2mm,minimum height=10mm,
    align=center,font=\small,inner sep=7pt},
  tealblock/.style={block,draw=Teal!40,fill=Tint},
  goldblock/.style={block,draw=Gold!40,fill=Warm},
  flow/.style={-{Latex[length=2.2mm]},draw=Ink!75,line width=.8pt},
  faintflow/.style={-{Latex[length=2mm]},draw=Muted,line width=.7pt}}
\begin{document}
\color{Ink}

% =====================================================================
% 封面
% =====================================================================
\begin{titlepage}
\thispagestyle{empty}
\vspace*{9mm}
{\sffamily\small\color{Teal}\bfseries ARCHITECTURE MANUAL\quad /\quad 01}\par
\vspace{14mm}
{\sffamily\fontsize{40}{46}\selectfont\bfseries YingluoNet}\par
\vspace{5mm}
{\sffamily\fontsize{28}{36}\selectfont\bfseries 结构说明书}\par
\vspace{5mm}
{\Large\color{Muted}v1.1\quad 从一个神经元到完整前向传播}\par
\vspace{11mm}
{\color{Teal}\rule{\linewidth}{1.2pt}}\par
\vspace{8mm}
\begin{minipage}{0.94\linewidth}
\large
每一层先生成相位，让相位之间发生一次交互；\par
再把正弦分支与 Tanh 分支按比例混合，\par
最后叠加到输入上，形成这一层的输出。
\end{minipage}\par
\vspace{12mm}
\begin{center}
\begin{tikzpicture}[node distance=6mm]
\node[tealblock,text width=35mm,minimum height=21mm] (p) {\textbf{相位耦合}\\[-1pt]\small 通过相位差计算修正};
\node[tealblock,text width=35mm,minimum height=21mm,right=7mm of p] (g) {\textbf{双分支门控}\\[-1pt]\small 混合 Sin 与 Tanh};
\node[tealblock,text width=35mm,minimum height=21mm,right=7mm of g] (r) {\textbf{残差连接}\\[-1pt]\small 保留输入并加入修正};
\draw[flow] (p)--(g);\draw[flow] (g)--(r);
\end{tikzpicture}
\end{center}
\vspace{7mm}
\begin{explain}{本书的阅读重点}
先看“数据经过哪些步骤”，再看“每一步为什么这样写”。公式均配有符号解释；正文提供一个可手算的两神经元例子，并在后半部分给出批量计算、参数初始化和梯度传播的对应关系。
\end{explain}
\vfill
{\small\color{Muted}核心定义：YingluoLayer v1.1\qquad 文档日期：2026 年 11 月 11 日}\par
\vspace{2mm}
{\small\color{Muted}数学结构\quad /\quad 前向传播\quad /\quad 参数说明\quad /\quad 接口说明}
\end{titlepage}
\setcounter{page}{2}

% =====================================================================
% 目录
% =====================================================================
\section*{阅读指南}
\begin{explain}{先记住一条主线}
\textbf{相位生成 $\rightarrow$ 相位交互 $\rightarrow$ 双分支混合 $\rightarrow$ 残差输出。}\\
整张网络，就是在输入与输出之间串联若干个这样的核心层。
\end{explain}
初次阅读可以顺着第 1--7 节走，重点看第 7 节的数值例子。需要与程序核对时，再看第 8--10 节；第 11--12 节用于理解训练与结构性质。最后一页可单独用作公式速查。
\vspace{2mm}
\begingroup
\renewcommand{\contentsname}{目录}
\tableofcontents
\endgroup
\vfill
{\small\color{Muted}\textbf{版本口径}\quad 本说明书依据本次对话中确定的 YingluoNet v1.1 公式与前向传播结构整理。核心层采用有向、零对角线耦合和 $1/m$ 缩放。输入编码、输出映射与训练损失分别说明，不混入核心层定义。数值算例为教学构造，不是训练权重。}

% =====================================================================
\newsection{先看全貌：输入、核心层、输出}
\subsection*{1.1\quad 三个部分各自做什么}
YingluoNet 可以拆成三个部分：\textbf{输入接口}把数据整理成向量；\textbf{核心层}反复变换特征；\textbf{输出接口}把最终特征读成所需的数值。核心层的名称是 \textbf{YingluoLayer}。
\begin{center}
\begin{tikzpicture}[node distance=5mm]
\node[block,text width=118mm] (x) {原始输入 $\boldsymbol{x}$};
\node[block,text width=118mm,below=of x] (enc) {输入编码 $\boldsymbol{z}=\operatorname{Encode}(\boldsymbol{x})\in\R^d$};
\node[tealblock,text width=118mm,below=of enc] (proj) {线性投影 $\boldsymbol{h}^{(0)}=W_{\rm in}\boldsymbol{z}+\boldsymbol{b}_{\rm in}\in\R^m$};
\node[tealblock,text width=118mm,below=of proj,minimum height=22mm] (core) {\textbf{连续串联 $L$ 个 YingluoLayer}\\
$\boldsymbol{h}^{(\ell+1)}=\YLayer_\ell(\boldsymbol{h}^{(\ell)})$，每层均为 $m\rightarrow m$};
\node[block,text width=118mm,below=of core] (read) {线性读出 $\boldsymbol{o}=W_{\rm out}\boldsymbol{h}^{(L)}+\boldsymbol{b}_{\rm out}\in\R^r$};
\node[goldblock,text width=118mm,below=of read] (out) {输出映射 $\widehat{\boldsymbol{y}}=T(\boldsymbol{o})$};
\draw[flow] (x)--(enc);\draw[flow] (enc)--(proj);\draw[flow] (proj)--(core);\draw[flow] (core)--(read);\draw[flow] (read)--(out);
\end{tikzpicture}
\end{center}
\subsection*{1.2\quad 核心层保持宽度不变}
这里 $d$ 是编码后的输入维度，$m$ 是隐藏宽度，$L$ 是核心层数量，$r$ 是输出维度。每层都接收一个 $m$ 维向量，再输出一个 $m$ 维向量，所以原输入与非线性修正可以直接相加。
\begin{explain}{“结构相同”与“参数相同”是两回事}
不同层采用相同的计算规则，但分别拥有自己的 $W_p,W_a,W_g,K$、偏置和 $\alpha$。它们不是同一组权重的重复使用。一次前向传播按层顺序计算，参数在这次计算中保持固定。
\end{explain}

% =====================================================================
\newsection{符号与维度：先分清参数和中间量}
正文先使用\textbf{单个样本的列向量}约定。向量用粗体表示，矩阵用大写字母表示；第 9 节再转换到程序常用的“每行一个样本”约定。
\subsection*{2.1\quad 这些量随当前输入变化}
\par\smallskip\noindent\begin{tabularx}{\linewidth}{@{}p{22mm}p{22mm}Y@{}}
\toprule
\textbf{符号}&\textbf{维度}&\textbf{通俗解释}\\\midrule
$\boldsymbol{h}$&$m\times1$&进入当前核心层的特征\\
$\boldsymbol{\theta}$&$m\times1$&每个神经元的原始相位\\
$\boldsymbol{c}$&$m\times1$&其他相位带来的修正\\
$\boldsymbol{\psi}$&$m\times1$&修正后的相位，$\boldsymbol{\psi}=\boldsymbol{\theta}+\boldsymbol{c}$\\
$\boldsymbol{p}$&$m\times1$&正弦分支的输出\\
$\boldsymbol{a}$&$m\times1$&Tanh 分支的输出\\
$\boldsymbol{g}$&$m\times1$&两个分支的混合比例\\
$\boldsymbol{f}$&$m\times1$&混合后的非线性特征\\
$\boldsymbol{h}_{+}$&$m\times1$&当前核心层的最终输出\\\bottomrule
\end{tabularx}\par\smallskip
\subsection*{2.2\quad 这些量通过训练调整}
\par\smallskip\noindent\begin{tabularx}{\linewidth}{@{}p{37mm}p{25mm}Y@{}}
\toprule
\textbf{参数}&\textbf{维度}&\textbf{作用}\\\midrule
$W_p,\boldsymbol{b}_p$&$m\times m,\ m$&生成原始相位\\
$W_a,\boldsymbol{b}_a$&$m\times m,\ m$&生成 Tanh 分支特征\\
$W_g,\boldsymbol{b}_g$&$m\times m,\ m$&生成门控比例\\
$K$&$m\times m$&控制相位之间的交互权重\\
$\alpha$&标量&控制非线性修正的层级强度\\\bottomrule
\end{tabularx}\par\smallskip
\subsection*{2.3\quad 三种运算不要混淆}
$W\boldsymbol{h}$ 是\textbf{矩阵乘法}，会混合不同坐标。$\boldsymbol{u}\odot\boldsymbol{v}$ 是\textbf{逐元素乘法}，即 $(\boldsymbol{u}\odot\boldsymbol{v})_i=u_iv_i$。$\sin(\boldsymbol{\theta})$ 等函数也是逐元素计算。
\begin{notebox}{最容易混淆的一点}
$W_g$ 和 $\boldsymbol{b}_g$ 是参数；$\boldsymbol{g}$ 是由当前输入算出来的结果。同样，$K$ 是参数，$\boldsymbol{c}$ 是由 $K$ 和当前相位共同算出的结果。$\alpha$ 是每层一个标量，不是每个神经元一个门控。
\end{notebox}

% =====================================================================
\newsection{相位生成与相位耦合}
\subsection*{3.1\quad 先给每个神经元一个相位}
\begin{equation}\label{eq:phase}
\boxed{\boldsymbol{\theta}=W_p\boldsymbol{h}+\boldsymbol{b}_p},\qquad
\theta_i=\sum_{k=1}^{m}W_{p,ik}h_k+b_{p,i}.
\end{equation}
这一步仍然是加权求和。“相位”指进入正弦、余弦运算的实数，按弧度解释；它不必是外部测得的物理相位，也不必限制在 $[0,2\pi)$ 内。
\subsection*{3.2\quad 再计算同层神经元的相互修正}
\begin{equation}\label{eq:coupling}
\boxed{c_i=\frac{1}{m}\sum_{j=1}^{m}K_{ij}\sin(\theta_j-\theta_i)},\qquad K_{ii}=0.
\end{equation}
\textbf{读法：}对第 $i$ 个神经元，依次比较其他神经元 $j$ 的相位；把每个相位差变成正弦信号，乘以对应的 $K_{ij}$，求和后除以宽度 $m$。
\par\smallskip\noindent\begin{tabularx}{\linewidth}{@{}p{27mm}p{24mm}Y@{}}
\toprule
相位差 $\theta_j-\theta_i$&正弦值&仅看该条连接的交互信号\\\midrule
$0$ 或 $\pi$&$0$&这一项为零\\
$\pi/2$&$1$&等于 $K_{ij}/m$\\
$-\pi/2$&$-1$&等于 $-K_{ij}/m$\\\bottomrule
\end{tabularx}\par\smallskip
\textbf{方向约定：}$K_{ij}$ 表示“$j$ 对 $i$ 的影响”。$K_{ij}$ 可以为正或负；$K_{ij}$ 与 $K_{ji}$ 独立，因而是有向耦合。对角线恒为零，因为 $\sin(\theta_i-\theta_i)=0$。
\purpose{用相位差显式构造特征之间的交互，而不是只对每个相位独立施加激活。$1/m$ 把交互写成按宽度缩放的加权和；它不代表每项相加后一定保持相同方差。}
\begin{notebox}{所有神经元同时更新}
先算出完整的原始相位 $\boldsymbol{\theta}$，再用这同一组相位算全部 $c_i$，最后得到 $\boldsymbol{\psi}=\boldsymbol{\theta}+\boldsymbol{c}$。这里没有“算完第一个就拿新值去算第二个”的顺序更新，也没有层内反复迭代。
\end{notebox}

% =====================================================================
\newsection{把成对交互改写成矩阵计算}
直接为每对神经元构造相位差，会产生一个 $m\times m$ 的中间矩阵。相同的计算可以只使用正弦向量、余弦向量和两次矩阵乘法。
\subsection*{4.1\quad 从三角恒等式开始}
\[
\sin(\theta_j-\theta_i)
=\sin\theta_j\cos\theta_i-\cos\theta_j\sin\theta_i.
\]
代入式~\eqref{eq:coupling}，把与 $j$ 无关的项提出求和号：
\begin{align}
c_i&=\frac{1}{m}\left[\cos\theta_i\sum_jK_{ij}\sin\theta_j
-\sin\theta_i\sum_jK_{ij}\cos\theta_j\right].
\end{align}
定义两个向量
\[
\boldsymbol{u}=\sin\boldsymbol{\theta},\qquad
\boldsymbol{v}=\cos\boldsymbol{\theta},
\]
于是得到与逐项求和完全相同的公式：
\begin{equation}\label{eq:vectorcoupling}
\boxed{\boldsymbol{c}=\frac{1}{m}\left[
\boldsymbol{v}\odot(K\boldsymbol{u})
-\boldsymbol{u}\odot(K\boldsymbol{v})\right].}
\end{equation}
\subsection*{4.2\quad 实际执行顺序}
\par\smallskip\noindent\begin{tabularx}{\linewidth}{@{}p{12mm}p{58mm}Y@{}}
\toprule
步骤&计算&得到什么\\\midrule
1&$\boldsymbol{u}=\sin\boldsymbol{\theta}$&全部相位的正弦\\
2&$\boldsymbol{v}=\cos\boldsymbol{\theta}$&全部相位的余弦\\
3&$\boldsymbol{r}=K\boldsymbol{u}$&正弦特征的加权汇总\\
4&$\boldsymbol{q}=K\boldsymbol{v}$&余弦特征的加权汇总\\
5&$\boldsymbol{c}=(\boldsymbol{v}\odot\boldsymbol{r}-\boldsymbol{u}\odot\boldsymbol{q})/m$&每个神经元的相位修正\\\bottomrule
\end{tabularx}\par\smallskip
\purpose{把显式的“逐对比较”改写成规则的张量运算。对一批 $b$ 个样本，主要运算量仍为 $O(bm^2)$，但不必显式保存 $b\times m\times m$ 的相位差张量。加上三条分支投影，每层矩阵运算主体约为 $5bm^2$ 次乘加，另有逐元素函数运算。}
\begin{notebox}{检查方向的简单方法}
若只有 $K_{12}$ 非零，则只有第 1 个神经元得到来自第 2 个神经元的直接相位修正。第 9 节采用行向量批量公式时，矩阵乘法会出现 $K^{\mathsf T}$，但方向约定完全不变。
\end{notebox}

% =====================================================================
\newsection{两条特征分支与一个门控}
\subsection*{5.1\quad 正弦分支：把修正后的相位读成特征}
\begin{equation}\label{eq:p}
\boldsymbol{\psi}=\boldsymbol{\theta}+\boldsymbol{c},\qquad
\boxed{\boldsymbol{p}=\sin\boldsymbol{\psi}.}
\end{equation}
这里的 $\sin$ 逐元素计算。相位修正发生在正弦函数\textbf{里面}，因此是先改相位、再产生输出；不是在正弦输出之后再额外加上 $\boldsymbol{c}$。
\subsection*{5.2\quad Tanh 分支：另一条独立的特征通道}
\begin{equation}\label{eq:a}
\boxed{\boldsymbol{a}=\tanh(W_a\boldsymbol{h}+\boldsymbol{b}_a).}
\end{equation}
这里 $\tanh z=(e^z-e^{-z})/(e^z+e^{-z})$，提供平滑、饱和的非周期响应。这一支直接读取 $\boldsymbol{h}$，不读取 $\boldsymbol{p}$ 或修正后的相位；它有独立的权重 $W_a$ 和偏置 $\boldsymbol{b}_a$。
\subsection*{5.3\quad 门控：按输入决定混合比例}
\begin{equation}\label{eq:g}
\boxed{\boldsymbol{g}=\sig(W_g\boldsymbol{h}+\boldsymbol{b}_g)},\qquad
\sig(z)=\frac{1}{1+e^{-z}}.
\end{equation}
每个 $g_i$ 都在 $(0,1)$ 之间，然后进行逐元素融合：
\begin{equation}\label{eq:f}
\boxed{\boldsymbol{f}=\boldsymbol{g}\odot\boldsymbol{p}
+(\boldsymbol{1}-\boldsymbol{g})\odot\boldsymbol{a}.}
\end{equation}
例如 $p_i=0.8$、$a_i=0.2$、$g_i=0.75$，则 $f_i=0.75\times0.8+0.25\times0.2=0.65$。每个神经元的比例分别计算，而且同一神经元面对不同输入时，比例也可以不同。
\begin{center}
\begin{tikzpicture}
\node[block,text width=20mm] (h) at (0,0) {$\boldsymbol{h}$};
\node[tealblock,text width=35mm] (p) at (-4.9,-1.8) {相位与耦合\\$\boldsymbol{p}=\sin(\boldsymbol{\theta}+\boldsymbol{c})$};
\node[tealblock,text width=35mm] (a) at (0,-1.8) {Tanh 分支\\$\boldsymbol{a}=\tanh(W_a\boldsymbol{h}+\boldsymbol{b}_a)$};
\node[goldblock,text width=35mm] (g) at (4.9,-1.8) {门控分支\\$\boldsymbol{g}=\sig(W_g\boldsymbol{h}+\boldsymbol{b}_g)$};
\node[block,text width=107mm] (f) at (0,-3.6) {$\boldsymbol{f}=\boldsymbol{g}\odot\boldsymbol{p}+(\boldsymbol{1}-\boldsymbol{g})\odot\boldsymbol{a}$};
\draw[flow] (h.south)--++(0,-.32)-|(p.north);
\draw[flow] (h)--(a);
\draw[flow] (h.south)--++(0,-.32)-|(g.north);
\draw[flow] (p.south)--++(0,-.35)-|([xshift=-30mm]f.north);
\draw[flow] (a)--(f);
\draw[flow] (g.south)--++(0,-.35)-|([xshift=30mm]f.north);
\end{tikzpicture}
\end{center}
\purpose{为同一层提供两种非线性响应，再让门控决定如何组合。$g_i$ 是分支的混合系数，不是该分支对最终预测的“贡献百分比”；后面还有残差和其他层。}

% =====================================================================
\newsection{残差融合：保留输入，再加入修正}
前面的分支计算得到了 $\boldsymbol{f}$，但它还不是当前层的最终输出。YingluoLayer 会把原输入 $\boldsymbol{h}$ 一起保留下来。
\begin{equation}\label{eq:res}
\boxed{\boldsymbol{h}_{+}=\frac{\boldsymbol{h}+\alpha\boldsymbol{f}}
{\sqrt{1+\alpha^2}}.}
\end{equation}
\subsection*{6.1\quad 两条路径怎样叠加}
把式~\eqref{eq:res} 拆开看：
\[
\boldsymbol{h}_{+}=
\underbrace{\frac{1}{\sqrt{1+\alpha^2}}}_{\text{输入路径系数}}\boldsymbol{h}
+\underbrace{\frac{\alpha}{\sqrt{1+\alpha^2}}}_{\text{修正路径系数}}\boldsymbol{f}.
\]
第一项沿着直接路径传递输入；第二项把非线性计算结果加进来。$\alpha$ 是\textbf{每层独立的一个可训练实数}，不是按输入生成的向量。
\par\smallskip\noindent\begin{tabularx}{\linewidth}{@{}p{25mm}p{37mm}p{37mm}Y@{}}
\toprule
$\alpha$&输入系数&修正系数&直观含义\\\midrule
$0$&$1$&$0$&该层为恒等映射\\
$0.1$&约 $0.9950$&约 $0.0995$&从较小的非线性修正开始\\
$1$&约 $0.7071$&约 $0.7071$&两条路径使用相同系数\\\bottomrule
\end{tabularx}\par\smallskip
\subsection*{6.2\quad 分母是尺度控制，不是标准化层}
两个路径系数的平方和等于 1，但系数之和通常不等于 1。因此，这不是两个向量的普通加权平均，也不是 LayerNorm 或 BatchNorm。
若把某一坐标在样本间的变化看作随机变量，则
\[
\operatorname{Var}(h_{+,i})=
\frac{\operatorname{Var}(h_i)+\alpha^2\operatorname{Var}(f_i)
+2\alpha\operatorname{Cov}(h_i,f_i)}{1+\alpha^2}.
\]
由于 $\boldsymbol{f}$ 本来就由 $\boldsymbol{h}$ 算出，两者一般有关联，不能把分母解释成“严格保持方差不变”。
\purpose{让输入到输出之间保留一条直接路径，并用可学习的 $\alpha$ 调节非线性修正。初始化采用 $\alpha=0.1$，使分支从小幅参与开始，而不是一开始就完全切断。}
\begin{notebox}{它与门控的分工}
$g_i$ 决定第 $i$ 个特征里 \textbf{Sin 与 Tanh 如何混合}；$\alpha$ 决定整层的\textbf{混合结果如何参与残差输出}。二者调节的位置不同。核心定义没有把 $\alpha$ 限制为正数或限制在 $[0,1]$。
\end{notebox}

% =====================================================================
\newsection{手算一次：两个神经元的完整例子}
下面取 $m=2$。所有数值都为便于理解而构造，三个偏置均为零。假设输入与参数如下：
\[
\boldsymbol{h}=\begin{bmatrix}0.5\\-0.5\end{bmatrix},\quad
W_p=\begin{bmatrix}1&1\\\pi/2&-\pi/2\end{bmatrix},\quad
K=\begin{bmatrix}0&1\\-1&0\end{bmatrix},\quad\alpha=0.1,
\]
\[
W_a=\begin{bmatrix}1&0\\0&1\end{bmatrix},\qquad
W_g=\begin{bmatrix}\ln3&-\ln3\\-\ln3&\ln3\end{bmatrix}.
\]
\subsection*{7.1\quad 先算相位与交互}
相位为 $\boldsymbol{\theta}=W_p\boldsymbol{h}=[0,\pi/2]^{\mathsf T}$。两项耦合分别为
\[
 c_1=\frac12(1)\sin(\pi/2-0)=0.5,\qquad
 c_2=\frac12(-1)\sin(0-\pi/2)=0.5.
\]
因此 $\boldsymbol{\psi}=[0.5,\pi/2+0.5]^{\mathsf T}$。注意第 2 项里，负权重乘上负正弦值，得到的修正仍为正数。
\subsection*{7.2\quad 再算两个分支、门控和残差}
Tanh 分支输入就是 $\boldsymbol{h}$。门控的加权和为 $[\ln3,-\ln3]^{\mathsf T}$，所以门控恰好是 $[0.75,0.25]^{\mathsf T}$。完整中间结果如下：
\begin{center}
\par\smallskip\noindent\begin{tabularx}{\linewidth}{@{}Yrr@{}}
\toprule
\textbf{中间量}&\textbf{神经元 1}&\textbf{神经元 2}\\\midrule
原始相位 $\theta_i$&$0$&$1.570796$\\
耦合修正 $c_i$&$0.500000$&$0.500000$\\
正弦输出 $p_i$&$0.479426$&$0.877583$\\
Tanh 输出 $a_i$&$0.462117$&$-0.462117$\\
门控 $g_i$&$0.750000$&$0.250000$\\
混合结果 $f_i$&$0.475098$&$-0.127192$\\
最终输出 $h_{+,i}$&$0.544793$&$-0.510175$\\\bottomrule
\end{tabularx}\par\smallskip
\end{center}
例如第 1 个神经元最后两步是：
\[
f_1=0.75(0.479426)+0.25(0.462117)\approx0.475098,
\]
\[
h_{+,1}=\frac{0.5+0.1(0.475098)}{\sqrt{1.01}}\approx0.544793.
\]
\begin{explain}{从这个例子记住什么}
耦合改变正弦里面的相位；门控混合两个分支；残差再把原输入加回来。隐藏特征可以为负，输出是否要求非负由网络末端的任务接口另行规定。
\end{explain}

% =====================================================================
\newsection{从核心层拼成完整网络}
\subsection*{8.1\quad 一般形式}
\begin{align}
\boldsymbol{z}&=\operatorname{Encode}(\boldsymbol{x}),&\boldsymbol{z}&\in\R^d,\notag\\
\boldsymbol{h}^{(0)}&=W_{\rm in}\boldsymbol{z}+\boldsymbol{b}_{\rm in},&\boldsymbol{h}^{(0)}&\in\R^m,\notag\\
\boldsymbol{h}^{(\ell+1)}&=\YLayer_\ell(\boldsymbol{h}^{(\ell)}),&\ell&=0,\ldots,L-1,\notag\\
\boldsymbol{o}&=W_{\rm out}\boldsymbol{h}^{(L)}+\boldsymbol{b}_{\rm out},&\boldsymbol{o}&\in\R^r,\notag\\
\widehat{\boldsymbol{y}}&=T(\boldsymbol{o}).&&\label{eq:whole}
\end{align}
$W_{\rm in}$ 的形状是 $m\times d$，$W_{\rm out}$ 的形状是 $r\times m$。中间各层的宽度保持为 $m$。
\subsection*{8.2\quad 离散输入接口：Embedding 与分类分数}
对于两个整数标签 $(a,b)$，设共享的可训练查找表为 $E\in\R^{p\times d_e}$。先查出两个向量，再按顺序拼接：
\[
\boldsymbol{z}=[E[a];E[b]]\in\R^{2d_e}.
\]
例如 $p=97$、$d_e=128$、$m=178$、$L=3$ 时，维度路线是
\[
(a,b)\ \longrightarrow\ 256\ \longrightarrow\ 178
\ \xrightarrow{\ 3\text{ 个核心层}\ }\ 178\ \longrightarrow\ 97.
\]
末端的 $\boldsymbol{o}$ 是 97 个实数分数（logits）。要选一个类别，可计算 $\arg\max_k o_k$；要表示类别概率，可计算 $e^{o_k}/\sum_j e^{o_j}$。这些是输出接口，不改变核心层。
\subsection*{8.3\quad 连续输入接口：时间与非负状态量}
对于时间 $\tau\in[\tau_{\min},\tau_{\max}]$，可以采用线性归一化
\[
 z=2\frac{\tau-\tau_{\min}}{\tau_{\max}-\tau_{\min}}-1.
\]
若末端有两个原始输出 $o_1,o_2$，并要求它们表示非负状态，可定义
\[
\widehat n=\spplus(o_1),\quad \widehat s=\spplus(o_2),\qquad
\spplus(v)=\ln(1+e^v).
\]
\begin{notebox}{结构层级要分开}
Embedding、时间归一化、Softplus 和分类读出都属于输入或输出接口。YingluoLayer 内部仍采用同一套相位、耦合、门控和残差公式。这里的 Softplus 是输出映射，不是把监督目标改成对数。
\end{notebox}

% =====================================================================
\newsection{批量前向传播：与张量维度对应}
程序通常把一个样本放在一行。为避免与偏置 $\boldsymbol{b}$ 混淆，下面用 $b$ 表示批量大小。输入矩阵为 $H\in\R^{b\times m}$。
\subsection*{9.1\quad 一次计算整批样本}
设 $\boldsymbol{1}_b$ 是长度为 $b$ 的全 1 列向量，$\mathsf{M}=\boldsymbol{1}_{m\times m}-I_m$ 是对角线为零的掩码。定义有效耦合矩阵 $K_{\rm eff}=K\odot\mathsf{M}$。
\begin{align}
\Theta &= H W_p^{\mathsf T}+\boldsymbol{1}_b\boldsymbol{b}_p^{\mathsf T},\notag\\
U&=\sin\Theta,\qquad V=\cos\Theta,\notag\\
C&=\frac{V\odot(U K_{\rm eff}^{\mathsf T})-U\odot(V K_{\rm eff}^{\mathsf T})}{m},\notag\\
P&=\sin(\Theta+C),\notag\\
A&=\tanh(H W_a^{\mathsf T}+\boldsymbol{1}_b\boldsymbol{b}_a^{\mathsf T}),\notag\\
G&=\sig(H W_g^{\mathsf T}+\boldsymbol{1}_b\boldsymbol{b}_g^{\mathsf T}),\notag\\
F&=G\odot P+(\boldsymbol{1}_{b\times m}-G)\odot A,\notag\\
H_+&=\frac{H+\alpha F}{\sqrt{1+\alpha^2}}.\label{eq:batch}
\end{align}
这些中间矩阵的形状全部为 $b\times m$。$\boldsymbol{1}_b\boldsymbol{b}_p^{\mathsf T}$ 表示把同一条偏置复制到每个样本；程序中通常通过广播完成。
\subsection*{9.2\quad 转置为什么出现}
单样本列向量写法是 $W_p\boldsymbol{h}$。把样本改成行向量后，同一个运算写成 $\boldsymbol{h}^{\mathsf T}W_p^{\mathsf T}$；堆叠多行就得到 $HW_p^{\mathsf T}$。耦合矩阵的转置同理。
\par\smallskip\noindent\begin{tabularx}{\linewidth}{@{}p{53mm}p{39mm}Y@{}}
\toprule
运算&形状检查&输出形状\\\midrule
$H W_p^{\mathsf T}$&$(b\times m)(m\times m)$&$b\times m$\\
$U K_{\rm eff}^{\mathsf T}$&$(b\times m)(m\times m)$&$b\times m$\\
$G\odot P$&两者均为 $b\times m$&$b\times m$\\
$H+\alpha F$&同形状相加&$b\times m$\\\bottomrule
\end{tabularx}\par\smallskip
\begin{explain}{批量不会引入样本间的耦合}
第 $r$ 行的相位修正只由第 $r$ 行的相位计算；$K$ 混合的是同一样本内部的神经元维度，不是不同样本。当前核心层没有跨样本的统计操作。
\end{explain}

% =====================================================================
\newsection{参数初始化与参数量}
\subsection*{10.1\quad 当前核心层的初始化口径}
\par\smallskip\noindent\begin{tabularx}{\linewidth}{@{}p{41mm}p{44mm}Y@{}}
\toprule
参数&初始设置&直观含义\\\midrule
$W_p,W_a,W_g$&Xavier 均匀初始化&三组矩阵分别随机生成\\
$\boldsymbol{b}_p,\boldsymbol{b}_a,\boldsymbol{b}_g$&全零&起点不加额外偏置\\
$K$&全零&起点的相位修正为零\\
$\alpha$&$0.1$&非线性分支从小幅参与开始\\
有效耦合对角线&始终为零&通过固定掩码屏蔽\\\bottomrule
\end{tabularx}\par\smallskip
这里对三组方阵采用增益为 1 的 Xavier 均匀初始化，即
\[
W_{ij}\sim\mathcal U\!\left(-\sqrt{\frac{6}{m+m}},\sqrt{\frac{6}{m+m}}\right)
=\mathcal U\!\left(-\sqrt{\frac{3}{m}},\sqrt{\frac{3}{m}}\right).
\]
输入层、输出层和 Embedding 的初始化属于外围接口配置。初始化只是参数的起点，训练后参数通常会改变。
\subsection*{10.2\quad 一层包含多少参数}
三条线性分支合计 $3m^2+3m$ 个参数；完整存储的耦合矩阵有 $m^2$ 个参数；另有一个 $\alpha$，所以
\begin{equation}
\boxed{P_{\rm layer}=4m^2+3m+1.}
\end{equation}
若只数直接参与前向计算的参数元素，耦合矩阵剩 $m(m-1)$ 项，因此 $P_{\rm layer,eff}=4m^2+2m+1$。常见的参数计数仍会包含被屏蔽的 $m$ 个对角元素；掩码本身不是可训练参数。
\subsection*{10.3\quad 完整网络的参数量}
编码输出维度为 $d$、输出维度为 $r$ 时，不含编码器的总量为
\begin{equation}
\boxed{P=(d+1)m+L(4m^2+3m+1)+(m+1)r.}
\end{equation}
若采用 $p\times d_e$ 的 Embedding 表，还需加上 $pd_e$。
\par\smallskip\noindent\begin{tabularx}{\linewidth}{@{}Yrr@{}}
\toprule
结构配置&完整存储计数&去除屏蔽对角线后\\\midrule
$d=1,m=98,L=3,r=2$&116,527&116,233\\
$d=256,m=178,L=3,r=97$，\newline 另加 $97\times128$ Embedding&457,338&456,804\\\bottomrule
\end{tabularx}\par\smallskip

% =====================================================================
\newsection{拓展阅读：梯度怎样穿过核心层}
这一节只解释局部梯度。设整个任务的损失为 $\mathcal L$，从后续网络传来的梯度是 $\boldsymbol{u}=\partial\mathcal L/\partial\boldsymbol{h}_{+}$。
\subsection*{11.1\quad 从残差输出进入混合特征}
\[
\boxed{\boldsymbol{q}=\frac{\partial\mathcal L}{\partial\boldsymbol{f}}
=\frac{\alpha}{\sqrt{1+\alpha^2}}\boldsymbol{u}.}
\]
之后沿两条分支分别传递：
\[
\frac{\partial\mathcal L}{\partial\boldsymbol{p}}
=\boldsymbol{q}\odot\boldsymbol{g},\qquad
\frac{\partial\mathcal L}{\partial\boldsymbol{a}}
=\boldsymbol{q}\odot(\boldsymbol{1}-\boldsymbol{g}).
\]
设门控的线性输入为 $\boldsymbol{v}_g=W_g\boldsymbol{h}+\boldsymbol{b}_g$，则
\[
\frac{\partial\mathcal L}{\partial\boldsymbol{v}_g}
=\boldsymbol{q}\odot(\boldsymbol{p}-\boldsymbol{a})
\odot\boldsymbol{g}\odot(\boldsymbol{1}-\boldsymbol{g}).
\]
这表明门控的梯度同时取决于两个分支的差别、当前门控值，以及后面传来的任务误差信号。
\subsection*{11.2\quad 相位耦合矩阵怎样得到梯度}
先穿过外层正弦，记
\[
\boldsymbol{\delta}=\frac{\partial\mathcal L}{\partial\boldsymbol{\psi}}
=\boldsymbol{q}\odot\boldsymbol{g}\odot\cos\boldsymbol{\psi}.
\]
对一条有效连接 $i\ne j$，有
\begin{equation}
\boxed{\frac{\partial\mathcal L}{\partial K_{ij}}
=\frac{\delta_i}{m}\sin(\theta_j-\theta_i).}
\end{equation}
因此，\textbf{$K=0$ 并不会自动让 $K$ 的梯度为零}。只要上游梯度和相位差提供非零信号，零初始化的耦合权重仍可改变。批量训练时，再按照损失的求和或平均方式汇总各样本贡献。
\subsection*{11.3\quad 为什么保留残差路径}
记 $J_f=\partial\boldsymbol{f}/\partial\boldsymbol{h}$，则当前层对输入的雅可比矩阵为
\begin{equation}
\boxed{J_{\YLayer}=\frac{I_m+\alpha J_f}{\sqrt{1+\alpha^2}}.}
\end{equation}
其中 $I_m$ 来自直接路径。$\alpha$ 很小时，这条路径接近恒等传播；但完整梯度仍取决于 $J_f$，不能仅凭残差形式就断言梯度尺度一定稳定。
\begin{notebox}{为什么不把 $\alpha$ 也初始化为零}
若 $\alpha=0$，该层输出恰好等于输入，但上面的 $\boldsymbol{q}$ 也为零，分支内部参数在这次反向传播中得不到来自该分支的梯度。取 $0.1$ 是让直接路径和分支训练同时保留的起点设置。
\end{notebox}

% =====================================================================
\newsection{读公式时的几个关键细节}
\subsection*{12.1\quad 这是一次前馈变换，而不是跨样本记忆}
对当前输入生成 $\boldsymbol{\theta}$，计算一次 $\boldsymbol{c}$，再输出 $\boldsymbol{h}_{+}$。每次前向传播都从当前输入重新计算中间量。数学定义中没有把上一次调用的相位保存到下一次调用。
\subsection*{12.2\quad 相位平移对耦合和输出的影响不同}
把所有相位加上同一个实数 $\delta$，相位差不变：
\[
(\theta_j+\delta)-(\theta_i+\delta)=\theta_j-\theta_i.
\]
所以 $\boldsymbol{c}$ 不变，但 $\sin(\boldsymbol{\theta}+\boldsymbol{c}+\delta\boldsymbol{1})$ 一般会变。\textbf{耦合项的共同平移不变性，不等于整个正弦分支的输出不变。}
\subsection*{12.3\quad 关闭一个机制后，公式还剩什么}
当 $K=0$ 时，核心层的混合特征为
\[
\boldsymbol{f}=\boldsymbol{g}\odot\sin(W_p\boldsymbol{h}+\boldsymbol{b}_p)
+(\boldsymbol{1}-\boldsymbol{g})\odot\tanh(W_a\boldsymbol{h}+\boldsymbol{b}_a).
\]
它仍然包含两条非线性分支、门控和残差。只有进一步让门控接近零，才会接近单独的 Tanh 分支；让门控接近一，则接近正弦分支。若 $\alpha=0$，整个核心层直接变为恒等映射。
\subsection*{12.4\quad 有界的是分支特征，不是任意输入下的完整输出}
因为 $|p_i|\le1$、$|a_i|<1$，且 $0<g_i<1$，所以
\[
|f_i|\le1,\qquad
|h_{+,i}|\le\frac{|h_i|+|\alpha|}{\sqrt{1+\alpha^2}}.
\]
后一个上界仍然依赖输入 $h_i$。这说明非线性修正的幅度受到限制，但没有把整张网络的输出固定限制在 $[-1,1]$。
\subsection*{12.5\quad 门控、耦合、损失：分别属于哪一层概念}
$K$ 决定相位如何交互；$\boldsymbol{g}$ 决定两个分支如何混合；$\alpha$ 决定混合特征如何加入残差。损失函数则放在完整网络输出之后，告诉优化过程“哪些输出更符合目标”。
\begin{explain}{最简洁的结构描述}
\textbf{YingluoNet v1.1 是把相位差交互、Sin/Tanh 门控混合与缩放残差组合起来的前馈网络。}理解它时，先分清三种作用发生的位置，再讨论具体的输入接口与输出接口。
\end{explain}

% =====================================================================
\clearpage
\appendix
\section{一页速查：核心公式与计算顺序}
\begin{explain}{单个样本 · 列向量约定}
输入 $\boldsymbol{h}\in\R^m$；$W_p,W_a,W_g,K\in\R^{m\times m}$；三个偏置均为 $m$ 维；$\alpha\in\R$；有效耦合矩阵满足 $K_{ii}=0$。
\end{explain}
\begin{align*}
\text{1. 原始相位}\qquad
&\boldsymbol{\theta}=W_p\boldsymbol{h}+\boldsymbol{b}_p,\\[5pt]
\text{2. 相位修正}\qquad
&c_i=\frac1m\sum_jK_{ij}\sin(\theta_j-\theta_i),\\[5pt]
\text{3. 正弦分支}\qquad
&\boldsymbol{p}=\sin(\boldsymbol{\theta}+\boldsymbol{c}),\\[5pt]
\text{4. Tanh 分支}\qquad
&\boldsymbol{a}=\tanh(W_a\boldsymbol{h}+\boldsymbol{b}_a),\\[5pt]
\text{5. 门控比例}\qquad
&\boldsymbol{g}=\sig(W_g\boldsymbol{h}+\boldsymbol{b}_g),\\[5pt]
\text{6. 特征混合}\qquad
&\boldsymbol{f}=\boldsymbol{g}\odot\boldsymbol{p}+(\boldsymbol{1}-\boldsymbol{g})\odot\boldsymbol{a},\\[5pt]
\text{7. 残差输出}\qquad
&\boxed{\boldsymbol{h}_{+}=\frac{\boldsymbol{h}+\alpha\boldsymbol{f}}{\sqrt{1+\alpha^2}}.}
\end{align*}
\subsection*{耦合的等价矩阵公式}
\[
\boxed{\boldsymbol{c}=\frac{
\cos\boldsymbol{\theta}\odot(K\sin\boldsymbol{\theta})
-\sin\boldsymbol{\theta}\odot(K\cos\boldsymbol{\theta})}{m}.}
\]
\subsection*{完整网络的连接关系}
\[
\boldsymbol{x}\xrightarrow{\mathrm{Encode}}\boldsymbol{z}
\xrightarrow{\mathrm{Linear}}\boldsymbol{h}^{(0)}
\xrightarrow{\YLayer_0}\cdots\xrightarrow{\YLayer_{L-1}}\boldsymbol{h}^{(L)}
\xrightarrow{\mathrm{Linear}}\boldsymbol{o}\xrightarrow{T}\widehat{\boldsymbol{y}}.
\]
\par\smallskip\noindent\begin{tabularx}{\linewidth}{@{}p{28mm}Y@{}}
\toprule
机制&一句话记忆\\\midrule
相位耦合 $K$&先看相位差，再改正弦里面的相位\\
门控 $\boldsymbol{g}$&对每个输入、每个特征计算 Sin/Tanh 混合比例\\
残差 $\alpha$&保留输入，同时加入经过缩放的非线性修正\\\bottomrule
\end{tabularx}\par\smallskip
\vfill
{\footnotesize\color{Muted}\textbf{源文件编译}\quad 本说明书为 UTF-8 编码的独立 \LaTeX{} 文件，使用 \texttt{ctexart}、\texttt{amsmath}、\texttt{TikZ} 等宏包；所有图均为矢量结构图。使用 XeLaTeX 编译两次即可更新目录与页码。}
\end{document}